%%% BEGIN fig:coll-transversity
\begin{figure}[htbp]
\centering
\begin{tikzpicture}[fd,
  fb lab/.style={font=\small},
  fb note/.style={font=\footnotesize,inner sep=1.5pt},
  fb hadron/.style={draw=black!70,line width=1.5pt,line cap=butt},
  fb harc/.style={draw=black!75,line width=1.6pt,line cap=round},
  fb odd/.style={draw=red!60!black,line width=1.2pt},
  fb oddlab/.style={font=\scriptsize,text=red!55!black,inner sep=1pt}]
  % ======================= (a) inclusive DIS =======================
  \begin{scope}
    \draw[quark] (-0.75,0.3) -- (-0.75,1.55);
    \draw[quark=0.25] (-0.75,1.55) -- (0.75,1.55);
    \draw[quark] (0.75,1.55) -- (0.75,0.3);
    \draw[photon] (-1.45,2.2) -- (-0.75,1.55);
    \draw[photon] (0.75,1.55) -- (1.45,2.2);
    \node[vertex] at (-0.75,1.55) {}; \node[vertex] at (0.75,1.55) {};
    \draw[fb hadron] (-1.9,0) -- (1.9,0);
    \draw[jet blob] (0,0) ellipse[x radius=1.15,y radius=0.4];
    \draw[fb odd] (0,0) ellipse[x radius=1.15,y radius=0.4];
    \node[font=\small] at (-0.52,0) {$h_1$};
    \draw[cut] (0,-0.6) -- (0,2.4);
    \node[fb oddlab,anchor=north] at (0.55,-0.42) {手征奇};
    \node[fb note,anchor=north west,text=gray!30!black] at (-0.72,1.5) {$\gamma^\mu$};
    \node[fb note,anchor=north east,text=gray!30!black] at (0.72,1.5) {$\gamma^\nu$};
    \node[font=\huge,text=red!65!black] at (1.55,0.85) {$\times$};
    \node[fb note,anchor=north,align=center] at (0,-0.85) {(a)\ 单举 DIS：只有一条费米子线，\\硬部分保持手征，$h_1$ 不出现};
  \end{scope}
  % ======================= (b) transversely polarized Drell--Yan =======================
  \begin{scope}[xshift=5.15cm]
    \draw[quark] (-0.75,0.3) -- (-0.75,1.2);
    \draw[quark] (-0.75,1.2) -- (-0.75,2.1);
    \draw[quark] (0.75,2.1) -- (0.75,1.2);
    \draw[quark] (0.75,1.2) -- (0.75,0.3);
    \draw[photon] (-0.75,1.2) -- (-1.65,1.2);
    \draw[photon] (0.75,1.2) -- (1.65,1.2);
    \node[vertex] at (-0.75,1.2) {}; \node[vertex] at (0.75,1.2) {};
    \draw[fb hadron] (-1.9,0) -- (1.9,0);
    \draw[fb hadron] (-1.9,2.4) -- (1.9,2.4);
    \draw[jet blob] (0,0) ellipse[x radius=1.15,y radius=0.4];
    \draw[jetB blob] (0,2.4) ellipse[x radius=1.15,y radius=0.4];
    \draw[fb odd] (0,0) ellipse[x radius=1.15,y radius=0.4];
    \draw[fb odd] (0,2.4) ellipse[x radius=1.15,y radius=0.4];
    \node[font=\small] at (-0.52,0) {$h_1$};
    \node[font=\small] at (-0.52,2.4) {$\bar h_1$};
    \draw[cut] (0,-0.6) -- (0,3.0);
    \node[fb oddlab,anchor=north] at (0.55,-0.42) {手征奇};
    \node[fb oddlab,anchor=south] at (0.55,2.82) {手征奇};
    \node[fb note,anchor=south east] at (-1.67,1.22) {$\gamma^*$};
    \node[font=\Large,text=green!45!black] at (1.62,0.62) {$\checkmark$};
    \node[fb note,anchor=north,align=center] at (0,-0.85) {(b)\ 横向极化 Drell--Yan：\\$h_1\otimes\bar h_1$（两个手征奇因子）};
  \end{scope}
  % ======================= (c) SIDIS with the Collins function =======================
  \begin{scope}[xshift=10.3cm]
    \draw[quark] (-0.75,0.3) -- (-0.75,1.2);
    \draw[quark] (-0.75,1.2) -- (-0.75,2.1);
    \draw[quark] (0.75,2.1) -- (0.75,1.2);
    \draw[quark] (0.75,1.2) -- (0.75,0.3);
    \draw[photon] (-1.65,1.55) -- (-0.75,1.2);
    \draw[photon] (0.75,1.2) -- (1.65,1.55);
    \node[vertex] at (-0.75,1.2) {}; \node[vertex] at (0.75,1.2) {};
    \draw[fb hadron] (-1.9,0) -- (1.9,0);
    \draw[jet blob] (0,0) ellipse[x radius=1.15,y radius=0.4];
    \draw[jetB blob] (0,2.4) ellipse[x radius=1.15,y radius=0.4];
    \draw[fb odd] (0,0) ellipse[x radius=1.15,y radius=0.4];
    \draw[fb odd] (0,2.4) ellipse[x radius=1.15,y radius=0.4];
    \draw[fb harc] (-0.38,2.76) to[out=80,in=180] (0,3.25) to[out=0,in=100] (0.38,2.76);
    \node[font=\small] at (-0.52,0) {$h_1$};
    \node[font=\small] at (-0.52,2.4) {$H_1^\perp$};
    \draw[cut] (0,-0.6) -- (0,3.5);
    \node[fb oddlab,anchor=north] at (0.55,-0.42) {手征奇};
    \node[fb oddlab,anchor=east] at (-1.2,2.4) {手征奇};
    \node[fb note,anchor=south west] at (0.08,3.27) {$h$};
    \node[font=\Large,text=green!45!black] at (1.62,0.62) {$\checkmark$};
    \node[fb note,anchor=north,align=center] at (0,-0.85) {(c)\ SIDIS（Collins 效应，TMD）：\\$h_1\otimes H_1^\perp$};
  \end{scope}
\end{tikzpicture}
\caption{横度分布 $h_1$ 的手征结构（割图示意；深红色边框标出手征奇的分布或碎裂函数）。光子顶点 $\gamma^\mu$ 保持（无质量）夸克的手征，手征奇的因子只能成对出现。
(a)~单举 DIS 只有一条经过硬部分的费米子线，$h_1$ 在领头扭度不出现（翻转手征需要夸克质量，被幂次压低）；
(b)~横向极化 Drell--Yan 中两个手征奇的分布相乘，$h_1\otimes\bar h_1$；(c)~SIDIS 中 $h_1$ 与手征奇的 Collins 碎裂函数 $H_1^\perp$ 相乘（TMD 观测量，见\cref{ch:tmd}）。横度分布的定义见\cref{def:coll-transversity}。}
\label{fig:coll-transversity}
\end{figure}
%%% END fig:coll-transversity
